% ID: SEL-D0018-BA-E02
% ESTADO: cotejada. Segunda pasada visual completa y borrador compilado revisados.
% Fuente: Selectividad/matematicas_aplicadas/matematicas_aplicadas_P_F.pdf
% SHA-256: fe6fd5b48407569d41169c80d71e6997050d6d06d1afd41cd639a553f9b712b5
% Página: 1; bloque A; ejercicio original 2.

Se considera la región del plano definida por las siguientes inecuaciones:
\[
x+y\leq 4\qquad x-y\geq -2\qquad x+3y\geq 2\qquad y\leq 2
\]
\begin{enumerate}[label=\alph*)]
\item \textbf{(1.5 puntos)} Represéntela gráficamente y determine sus vértices.
\item \textbf{(0.25 puntos)} Indique razonadamente si el punto $(4,-0.75)$ pertenece a dicha región.
\item \textbf{(0.75 puntos)} \textquestiondown{}En qué puntos de la región anterior la función $F(x,y)=x+y$ alcanza los valores máximo y mínimo y cuáles son estos valores?
\end{enumerate}
